\section{堆、Top K 和贪心章正文案例推导补强}

这一节补齐本章正文出现过的 LeetCode 案例推导。阅读顺序统一按“题目说明、样例入口、暴力基线、暴力过程、重复工作、优化条件、相邻状态、状态复用、指针和字段、代码落点、正确性、边界实验、复杂度变化、迁移追问”展开；目的不是新增题量，而是把正文中的案例都讲成可复述、可证明、可迁移的解题链。

\subsection{LeetCode 347 前 K 个高频元素}

\textbf{题目说明。} LeetCode 347 前 K 个高频元素 的题面入口要先还原成具体任务：先统计频次，再只维护前 K 个候选。

\textbf{样例入口。} 拿 [1, 1, 1, 2, 2, 3]、k=2 手算，频次是 1->3、2->2、3->1，前两个高频是 1 和 2。

\textbf{暴力基线。} 暴力先统计频次，再把所有不同元素按频次完整排序，取前 K 个。

\textbf{暴力过程。} 以 [1, 1, 1, 2, 2, 3]、k=2 为例，频次为 1->3、2->2、3->1，答案是 1 和 2。

\textbf{重复工作。} 题面模型：先统计频次，再只维护前 K 个候选。暴力版的慢点要落到本题：浪费在完整排序所有不同元素。只要前 K 高频，低于当前第 K 高频边界的元素可以丢弃。后面的优化只消掉这类重复工作，不能改变答案集合。

\textbf{优化条件。} 本题的关键观察是：小顶堆大小超过 K 就弹出最低频；桶排序利用频次上界。这句话必须能翻译成可维护状态；如果题目条件改变后观察不再成立，就不能继续套原来的写法。

\textbf{相邻状态。} 规律是先用哈希计数，再用大小为 k 的小顶堆或桶按频次取 Top K；频次相同不影响题目集合答案。把这个特例继续拆开看：堆顶必须正好代表下一步最需要处理的候选；若堆里可能有过期对象，读取堆顶前要先清理。堆只解决动态极值，不自动证明被弹出或被丢弃的元素无用。

\textbf{状态复用。} 先用哈希表计数。再用大小为 K 的小顶堆保存候选，堆顶是当前前 K 中频次最低者；超过 K 就弹出。手算时只改被新输入影响的那一小部分，而不是回头重算完整候选。

\textbf{指针和字段。} 状态字段要能说清：freq[value] 是频次，heap 元素是 (frequency, value)，heap.top 是淘汰边界。手算时按这个协议跟踪：插入频次 3 和 2 后堆满；频次 1 的元素入堆会被立刻弹出，留下 1 和 2。

\textbf{代码落点。} 关键是比较器按频次建小顶堆。题目通常不要求输出顺序；若要求顺序，最后还要按规则排序结果。关键代码行必须能逐句翻译成“旧状态怎样变成新状态”，否则就只是模板。

\textbf{正确性。} 堆始终保存已处理元素中频次最高的 K 个；任何被弹出的元素频次不高于堆中 K 个候选，不可能进入答案。

\textbf{边界实验。} 先用K 等于不同元素数、频次相同、结果顺序不要求做反例检查。实验时覆盖K 等于不同元素数、频次相同、结果顺序不要求，记录堆顶是否过期、比较器是否让正确候选在堆顶、K 或窗口大小为边界值时是否稳定。

\textbf{复杂度变化。} 完整排序 O(m log m)，m 为不同元素数；堆法 O(n+m log k)，空间 O(m+k)。

\textbf{迁移追问。} 如果题目改成“若要按频次和字典序排序，比较器要同时处理两个键”，先判断原状态是否仍然覆盖新答案集合，再决定是微调边界、改 value 语义，还是换算法。

\subsection{LeetCode 295 数据流的中位数}

\textbf{题目说明。} LeetCode 295 数据流的中位数 的题面入口要先还原成具体任务：数据流需要动态维护较小一半和较大一半。

\textbf{样例入口。} 拿数据流 1、2、3 手算，插入 1 后中位数 1；插入 2 后两半是 [1] 和 [2]，中位数 1.5；插入 3 后右半多一个，要平衡成左半 [2, 1]、右半 [3]。

\textbf{暴力基线。} 暴力每插入一个数就重新排序所有已到达元素，然后按长度奇偶取中位数。

\textbf{暴力过程。} 数据流 1、2、3 中，插入 1 中位数是 1；插入 2 后中位数是 (1+2)/2；插入 3 后中位数是 2。

\textbf{重复工作。} 题面模型：数据流需要动态维护较小一半和较大一半。暴力版的慢点要落到本题：浪费在每次都维护完整排序。中位数只需要较小一半的最大值和较大一半的最小值。后面的优化只消掉这类重复工作，不能改变答案集合。

\textbf{优化条件。} 本题的关键观察是：大顶堆保存较小一半，小顶堆保存较大一半，大小差不超过一。这句话必须能翻译成可维护状态；如果题目条件改变后观察不再成立，就不能继续套原来的写法。

\textbf{相邻状态。} 规律是大顶堆保存较小一半，小顶堆保存较大一半，大小差最多 1。把这个特例继续拆开看：堆顶必须正好代表下一步最需要处理的候选；若堆里可能有过期对象，读取堆顶前要先清理。堆只解决动态极值，不自动证明被弹出或被丢弃的元素无用。

\textbf{状态复用。} 大顶堆 small 保存较小一半，小顶堆 large 保存较大一半。保持 small.size() 与 large.size() 差不超过 1，且 small.top()<=large.top()。手算时只改被新输入影响的那一小部分，而不是回头重算完整候选。

\textbf{指针和字段。} 状态字段要能说清：small.top 是左半最大值，large.top 是右半最小值；奇数时较大堆顶是中位数，偶数时两个堆顶平均。手算时按这个协议跟踪：插入 1 放 small；插入 2 后移动平衡成 small=[1]、large=[2]；插入 3 后 large 多，再移动 2 到 small，中位数 2。

\textbf{代码落点。} 关键是先按顺序不变量放入堆，再重平衡大小。求平均时用 double，并防止两个 int 相加溢出。关键代码行必须能逐句翻译成“旧状态怎样变成新状态”，否则就只是模板。

\textbf{正确性。} 两个堆按大小分割数据流，中位数只可能位于左半最大或右半最小；大小不变量保证取堆顶即可。

\textbf{边界实验。} 先用偶数个元素、负数、重复值、平均值溢出做反例检查。实验时覆盖偶数个元素、负数、重复值、平均值溢出，记录堆顶是否过期、比较器是否让正确候选在堆顶、K 或窗口大小为边界值时是否稳定。

\textbf{复杂度变化。} 每次重新排序 O(n log n)；双堆 addNum O(log n)，findMedian O(1)，空间 O(n)。

\textbf{迁移追问。} 如果题目改成“若还要删除任意元素，需要延迟删除哈希表”，先判断原状态是否仍然覆盖新答案集合，再决定是微调边界、改 value 语义，还是换算法。

\subsection{LeetCode 23 合并 K 个升序链表}

\textbf{题目说明。} LeetCode 23 合并 K 个升序链表 的题面入口要先还原成具体任务：每条链表当前头节点都是可能的下一个最小元素。

\textbf{样例入口。} 拿三条链表头 1、1、2 手算，每次答案只需要当前最小头节点；弹出某条链的头后，只有它的后继会成为新候选。暴力每轮扫 k 个头，堆把“找当前最小头”变成对数操作。

\textbf{暴力基线。} 暴力每轮扫描 K 个链表头，找到当前最小头节点接到结果后推进该链。

\textbf{暴力过程。} 若三条链当前头是 1、1、2，第一步只需从这三个头里选最小；弹出某条链的 1 后，只有它的后继会成为新候选。

\textbf{重复工作。} 题面模型：每条链表当前头节点都是可能的下一个最小元素。暴力版的慢点要落到本题：浪费在每输出一个节点都线性扫 K 个头。候选集合动态变化，但每次只需要最小头节点，适合小顶堆。后面的优化只消掉这类重复工作，不能改变答案集合。

\textbf{优化条件。} 本题的关键观察是：小顶堆保存每条非空链的当前头，弹出最小节点后把它的后继入堆。这句话必须能翻译成可维护状态；如果题目条件改变后观察不再成立，就不能继续套原来的写法。

\textbf{相邻状态。} 规律是堆里始终最多保存每条非空链的一个当前头。把这个特例继续拆开看：堆顶必须正好代表下一步最需要处理的候选；若堆里可能有过期对象，读取堆顶前要先清理。堆只解决动态极值，不自动证明被弹出或被丢弃的元素无用。

\textbf{状态复用。} 小顶堆保存每条非空链的当前头节点。弹出最小节点接入结果；若该节点有 next，把 next 入堆。手算时只改被新输入影响的那一小部分，而不是回头重算完整候选。

\textbf{指针和字段。} 状态字段要能说清：heap 中元素是 ListNode*，比较依据是 node->val；tail 是结果尾；node->next 是该链的新候选。手算时按这个协议跟踪：头节点 1、1、2 入堆；弹出第一个 1 后把它的后继入堆，堆仍保存每条链当前最小暴露节点。

\textbf{代码落点。} 关键是堆里保存节点指针而不是复制值，这样可以复用原节点。比较器要写成小顶堆语义，C++ priority\_queue 默认是大顶堆。关键代码行必须能逐句翻译成“旧状态怎样变成新状态”，否则就只是模板。

\textbf{正确性。} 每条有序链未输出部分的最小值就是当前头；所有链当前头中的最小者就是全局下一个输出节点。

\textbf{边界实验。} 先用空链表数组、某条链为空、重复值、堆里保存节点指针还是值做反例检查。实验时覆盖空链表数组、某条链为空、重复值、堆里保存节点指针还是值，记录堆顶是否过期、比较器是否让正确候选在堆顶、K 或窗口大小为边界值时是否稳定。

\textbf{复杂度变化。} 每轮扫 K 个头是 O(NK)。堆中最多 K 个节点，每个节点入堆出堆一次，总时间 O(N log K)，空间 O(K)。

\textbf{迁移追问。} 如果题目改成“也可以分治两两合并，堆更适合 K 较大且链长不均的场景”，先判断原状态是否仍然覆盖新答案集合，再决定是微调边界、改 value 语义，还是换算法。

\subsection{LeetCode 55 跳跃游戏}

\textbf{题目说明。} LeetCode 55 跳跃游戏 的题面入口要先还原成具体任务：扫描过程中只关心当前能到达的最远位置。

\textbf{样例入口。} 拿 [2, 3, 1, 1, 4] 手算，下标 0 能覆盖到 2；走到下标 1 时最远可达更新到 4，已经覆盖终点。再拿 [0, 1]，扫描到下标 1 时已经超过最远可达，失败。

\textbf{暴力基线。} 暴力把每个下标能跳到的位置都当成分支搜索，尝试是否存在一条路径能到终点。

\textbf{暴力过程。} 以 [2, 3, 1, 1, 4] 为例，从 0 可以跳到 1 或 2；若到 1，最远能继续覆盖到 4。以 [0, 1] 为例，从 0 一步也走不出去。

\textbf{重复工作。} 题面模型：扫描过程中只关心当前能到达的最远位置。暴力版的慢点要落到本题：浪费在枚举具体路径。判断能否到达终点时，不关心走的是哪条路径，只关心目前所有可达位置能覆盖到的最远下标。后面的优化只消掉这类重复工作，不能改变答案集合。

\textbf{优化条件。} 本题的关键观察是：如果下标超过最远可达，任何之前路径都不能到达它；否则用它扩展边界。这句话必须能翻译成可维护状态；如果题目条件改变后观察不再成立，就不能继续套原来的写法。

\textbf{相邻状态。} 规律是只维护当前能到达的最远边界，任何超过边界的位置都不可能由前面路径到达。把这个特例继续拆开看：先构造一个非贪心选择，再比较两者留下的资源、右边界或剩余时间。只有能把非贪心第一步交换成贪心第一步且不变差，局部选择才是规律。

\textbf{状态复用。} 扫描下标 i，维护 farthest。若 i>farthest，说明当前下标不可达，直接失败；否则用 i+nums[i] 更新 farthest。手算时只改被新输入影响的那一小部分，而不是回头重算完整候选。

\textbf{指针和字段。} 状态字段要能说清：i 是当前检查位置，farthest 是所有已可达位置能跳到的最远边界，n-1 是目标边界。手算时按这个协议跟踪：样例中 i=0 后 farthest=2；i=1 可达并把 farthest 更新为 4，已经覆盖终点。

\textbf{代码落点。} 关键顺序是先判断当前位置是否超过 farthest；若没超过，再用 i+nums[i] 扩展最远可达位置。数组长度为 1 时起点就是终点。关键代码行必须能逐句翻译成“旧状态怎样变成新状态”，否则就只是模板。

\textbf{正确性。} 在扫描到 i 前，farthest 汇总了所有可达下标的最远覆盖；若 i 超过它，任何旧路径都到不了 i，更不可能通过 i 之后的位置到终点。

\textbf{边界实验。} 先用第一个元素为零、数组长度一、恰好到达最后、远超最后做反例检查。实验时除了第一个元素为零、数组长度一、恰好到达最后、远超最后，还要故意替换成一个看似合理但错误的局部规则，构造反例。能解释错误贪心为何失败，才算理解正确贪心。

\textbf{复杂度变化。} 路径搜索最坏指数级；贪心边界扫描 O(n) 时间、O(1) 空间。

\textbf{迁移追问。} 如果题目改成“求最少跳数时把可达区间看成 BFS 层”，先判断原状态是否仍然覆盖新答案集合，再决定是微调边界、改 value 语义，还是换算法。

\subsection{LeetCode 45 跳跃游戏 II}

\textbf{题目说明。} LeetCode 45 跳跃游戏 II 的题面入口要先还原成具体任务：每一次跳跃可以看成 BFS 的一层，当前层内所有位置共享同一次跳数。

\textbf{样例入口。} 拿 [2, 3, 1, 1, 4] 手算，起点这一层能到下标 1 和 2；扫描这一层时，下一层最远可到 4，走完当前层末尾才把步数加一。

\textbf{暴力基线。} 暴力 BFS 可以从每个当前下标扩展所有可跳位置，按层找第一次到终点的步数。

\textbf{暴力过程。} 以 [2, 3, 1, 1, 4] 为例，第一层是从 0 能到的 1、2；扫描这一层时，下一层最远可到 4，所以答案是 2。

\textbf{重复工作。} 题面模型：每一次跳跃可以看成 BFS 的一层，当前层内所有位置共享同一次跳数。暴力版的慢点要落到本题：浪费在显式保存一层内所有节点。数组下标的可达层是连续区间，只需要当前层右边界和下一层最远边界。后面的优化只消掉这类重复工作，不能改变答案集合。

\textbf{优化条件。} 本题的关键观察是：扫描当前层时维护下一层最远边界，到达当前层末尾才增加跳数。这句话必须能翻译成可维护状态；如果题目条件改变后观察不再成立，就不能继续套原来的写法。

\textbf{相邻状态。} 规律是把可达区间看成 BFS 层，当前层结束时切到下一层；不能在每个位置都加跳数。把这个特例继续拆开看：先构造一个非贪心选择，再比较两者留下的资源、右边界或剩余时间。只有能把非贪心第一步交换成贪心第一步且不变差，局部选择才是规律。

\textbf{状态复用。} current\_end 表示当前跳数能覆盖的最右位置，farthest 表示下一跳能覆盖的最远位置。扫描到 current\_end 时，步数加一并切换到下一层。手算时只改被新输入影响的那一小部分，而不是回头重算完整候选。

\textbf{指针和字段。} 状态字段要能说清：i 是扫描位置，jumps 是已完成层数，current\_end 是本层末尾，farthest 是下一层末尾。手算时按这个协议跟踪：i=0 时 farthest=2，走到 current\_end=0 后 jumps=1、current\_end=2；扫描 i=1 时 farthest 更新到 4，走到 i=2 后 jumps=2。

\textbf{代码落点。} 关键是循环通常只扫到 n-2，因为到达最后一个位置不需要再跳。不能在每个位置都 jumps++，只能在层结束时加。关键代码行必须能逐句翻译成“旧状态怎样变成新状态”，否则就只是模板。

\textbf{正确性。} 所有当前跳数可达的位置构成一个连续区间；在这个区间内选择能让下一层覆盖最远的位置，不会比保存全部 BFS 节点差。

\textbf{边界实验。} 先用数组长度一、刚好到达、存在多个可选跳点、不要在每个位置都加跳数做反例检查。实验时除了数组长度一、刚好到达、存在多个可选跳点、不要在每个位置都加跳数，还要故意替换成一个看似合理但错误的局部规则，构造反例。能解释错误贪心为何失败，才算理解正确贪心。

\textbf{复杂度变化。} 显式 BFS 分支多；区间层扫描 O(n) 时间、O(1) 空间。

\textbf{迁移追问。} 如果题目改成“若问能否到达，只需维护全局最远可达位置”，先判断原状态是否仍然覆盖新答案集合，再决定是微调边界、改 value 语义，还是换算法。

\subsection{LeetCode 621 任务调度器}

\textbf{题目说明。} LeetCode 621 任务调度器 的题面入口要先还原成具体任务：相同任务之间有冷却间隔，核心受最高频任务骨架限制。

\textbf{样例入口。} 拿 A, A, A, B, B, B、n=2 手算，若只按顺序执行 A 会违反冷却，必须在相同任务之间间隔两个时间单位。

\textbf{暴力基线。} 暴力按时间模拟，每个时刻扫描所有任务，找一个未冷却且剩余次数最多的任务执行，否则空闲。

\textbf{暴力过程。} 任务 A, A, A, B, B, B、n=2 中，相同 A 之间必须间隔两个单位；一种最短安排是 A B idle A B idle A B。

\textbf{重复工作。} 题面模型：相同任务之间有冷却间隔，核心受最高频任务骨架限制。暴力版的慢点要落到本题：浪费在逐时刻扫描所有任务。更重要的是最高频任务决定冷却骨架，其他任务只是填空。后面的优化只消掉这类重复工作，不能改变答案集合。

\textbf{优化条件。} 本题的关键观察是：可用大顶堆模拟，也可用最高频公式计算最短时间。这句话必须能翻译成可维护状态；如果题目条件改变后观察不再成立，就不能继续套原来的写法。

\textbf{相邻状态。} 规律是最高频任务决定骨架长度，也可用堆按剩余次数模拟，每轮安排最多 n+1 个不同任务。把这个特例继续拆开看：堆顶必须正好代表下一步最需要处理的候选；若堆里可能有过期对象，读取堆顶前要先清理。堆只解决动态极值，不自动证明被弹出或被丢弃的元素无用。

\textbf{状态复用。} 公式法统计最大频次 maxFreq 和拥有最大频次的任务数 maxCount。骨架长度为 (maxFreq-1)*(n+1)+maxCount，答案取它和任务总数的较大值。手算时只改被新输入影响的那一小部分，而不是回头重算完整候选。

\textbf{指针和字段。} 状态字段要能说清：maxFreq 是最高频任务次数，maxCount 是最高频任务种类数，n 是冷却间隔，tasks.size() 是没有空闲时的下界。手算时按这个协议跟踪：A 和 B 都出现 3 次，maxFreq=3、maxCount=2，骨架为 (3-1)*3+2=8。

\textbf{代码落点。} 关键是最后 return max(tasks.size(), skeleton)。若 n=0，任务总数就是答案。关键代码行必须能逐句翻译成“旧状态怎样变成新状态”，否则就只是模板。

\textbf{正确性。} 最高频任务之间至少有 maxFreq-1 个间隔块，每块长度 n+1；多个最高频任务占据最后一列。其他任务只能填充空位，不能突破任务总数下界。

\textbf{边界实验。} 先用冷却为零、多个最高频、空闲被其他任务填满做反例检查。实验时覆盖冷却为零、多个最高频、空闲被其他任务填满，记录堆顶是否过期、比较器是否让正确候选在堆顶、K 或窗口大小为边界值时是否稳定。

\textbf{复杂度变化。} 堆模拟 O(T log C)，公式法统计频次 O(T+C)，空间 O(C)，C 为任务种类数。

\textbf{迁移追问。} 如果题目改成“若任务有不同释放时间或执行时长，就变成更一般的调度问题”，先判断原状态是否仍然覆盖新答案集合，再决定是微调边界、改 value 语义，还是换算法。

\subsection{LeetCode 134 加油站}

\textbf{题目说明。} LeetCode 134 加油站 的题面入口要先还原成具体任务：绕行一圈是否可行取决于总油量差和局部剩余油量。

\textbf{样例入口。} 拿 gas=[1, 2, 3, 4, 5]、cost=[3, 4, 5, 1, 2] 手算，从 0 出发中途油量为负，0 到失败点之间任何站都不能作为起点；从 3 出发可以绕完。

\textbf{暴力基线。} 暴力把每个加油站都当起点，模拟绕一圈，若油量中途为负就换下一个起点。

\textbf{暴力过程。} gas=[1, 2, 3, 4, 5]、cost=[3, 4, 5, 1, 2] 中，从 0 出发会在中途油量为负；从 3 出发可以完成一圈。

\textbf{重复工作。} 题面模型：绕行一圈是否可行取决于总油量差和局部剩余油量。暴力版的慢点要落到本题：浪费在失败起点之后仍逐个重试中间站。若从 start 到 i 的累计油量已经为负，那么 start..i 中任何站都不可能作为可行起点。后面的优化只消掉这类重复工作，不能改变答案集合。

\textbf{优化条件。} 本题的关键观察是：总和为负必不可行；扫描时若当前剩余为负，起点只能放到下一站。这句话必须能翻译成可维护状态；如果题目条件改变后观察不再成立，就不能继续套原来的写法。

\textbf{相邻状态。} 规律是总油量差为负则无解，局部剩余为负时起点跳到下一站。把这个特例继续拆开看：先构造一个非贪心选择，再比较两者留下的资源、右边界或剩余时间。只有能把非贪心第一步交换成贪心第一步且不变差，局部选择才是规律。

\textbf{状态复用。} 扫描一次，total 记录总油量差，tank 记录当前候选起点到当前位置的剩余油量，start 是当前候选起点。手算时只改被新输入影响的那一小部分，而不是回头重算完整候选。

\textbf{指针和字段。} 状态字段要能说清：diff=gas[i]-cost[i]，tank 是局部剩余，total 是全局可行性，start 是下一次尝试起点。手算时按这个协议跟踪：样例前几站累计 tank 变负后，把 start 移到下一站并清零 tank；最后 total>=0，返回最终 start=3。

\textbf{代码落点。} 关键是 total 全程累加，tank<0 时 start=i+1、tank=0。循环结束后 total<0 返回 -1，否则返回 start。关键代码行必须能逐句翻译成“旧状态怎样变成新状态”，否则就只是模板。

\textbf{正确性。} tank<0 说明从当前 start 到 i 的总补给不足；其中任意中间站作为起点，少了前面一段补给，只会更差，故可整体跳过。

\textbf{边界实验。} 先用总和刚好为零、多个可行起点、某站油量为零、消耗大于补给做反例检查。实验时除了总和刚好为零、多个可行起点、某站油量为零、消耗大于补给，还要故意替换成一个看似合理但错误的局部规则，构造反例。能解释错误贪心为何失败，才算理解正确贪心。

\textbf{复杂度变化。} 暴力 O(\ensuremath{n^{2}})，一次扫描 O(n) 时间、O(1) 空间。

\textbf{迁移追问。} 如果题目改成“它的领先性证明是前一段任意站点都不能作为起点”，先判断原状态是否仍然覆盖新答案集合，再决定是微调边界、改 value 语义，还是换算法。

\subsection{LeetCode 763 划分字母区间}

\textbf{题目说明。} LeetCode 763 划分字母区间 的题面入口要先还原成具体任务：每个片段必须覆盖其中所有字符的最后出现位置。

\textbf{样例入口。} 拿 ababcbacadefegdehijhklij 手算，字符 a 的最后位置在 8，所以第一个片段至少到 8；扫描中 b、c 的最后位置也不超过 8，到 8 才能切。

\textbf{暴力基线。} 暴力尝试每个切分位置，并检查切出来的片段中是否有字符还会在后面再次出现。

\textbf{暴力过程。} 字符串 ababcbacadefegdehijhklij 中，a 最后出现在 8，所以第一个片段不可能早于 8 结束；扫描到 b、c 时也要把片段右端扩到它们的最后位置。

\textbf{重复工作。} 题面模型：每个片段必须覆盖其中所有字符的最后出现位置。暴力版的慢点要落到本题：浪费在每次试切后重新检查片段字符。每个字符的最后出现位置可以预处理，扫描时维护当前片段必须覆盖的最远右端。后面的优化只消掉这类重复工作，不能改变答案集合。

\textbf{优化条件。} 本题的关键观察是：扫描维护当前片段最远右端，到达右端即可切分。这句话必须能翻译成可维护状态；如果题目条件改变后观察不再成立，就不能继续套原来的写法。

\textbf{相邻状态。} 规律是当前片段右端是片段内所有字符最后出现位置的最大值。把这个特例继续拆开看：先构造一个非贪心选择，再比较两者留下的资源、右边界或剩余时间。只有能把非贪心第一步交换成贪心第一步且不变差，局部选择才是规律。

\textbf{状态复用。} last[ch] 保存字符最后出现下标。扫描字符串时，right=max(right, last[s[i]])；当 i==right，当前片段可以切开。手算时只改被新输入影响的那一小部分，而不是回头重算完整候选。

\textbf{指针和字段。} 状态字段要能说清：start 是当前片段起点，i 是扫描位置，right 是当前片段内所有字符最后位置的最大值。手算时按这个协议跟踪：第一个片段从 0 开始，right 先因 a 变成 8；扫描到 i=8 时，片段内所有字符都不会再出现，可以切出长度 9。

\textbf{代码落点。} 关键是先预处理 last，再扫描更新 right。切分后记录 right-start+1，并令 start=i+1。关键代码行必须能逐句翻译成“旧状态怎样变成新状态”，否则就只是模板。

\textbf{正确性。} 片段必须覆盖内部每个字符的最后出现位置，所以 right 之前不能切；当 i==right 时，片段内字符都不会越界，立即切能最大化片段数量。

\textbf{边界实验。} 先用单字符、所有字符相同、字符多次跨段做反例检查。实验时除了单字符、所有字符相同、字符多次跨段，还要故意替换成一个看似合理但错误的局部规则，构造反例。能解释错误贪心为何失败，才算理解正确贪心。

\textbf{复杂度变化。} 暴力试切最坏 O(\ensuremath{n^{2}})，预处理加扫描 O(n) 时间，字符集空间 O(1) 或 O(字符集)。

\textbf{迁移追问。} 如果题目改成“这是最早合法切分贪心，能最大化片段数量”，先判断原状态是否仍然覆盖新答案集合，再决定是微调边界、改 value 语义，还是换算法。

\subsection{LeetCode 502 IPO}

\textbf{题目说明。} LeetCode 502 IPO 的题面入口要先还原成具体任务：当前资本决定可解锁项目集合，目标是在最多 k 次选择中最大化最终资本。

\textbf{样例入口。} 拿 k=2、w=0、profits=[1, 2, 3]、capital=[0, 1, 1] 手算，初始只能做资本 0 的项目，做完利润 1 后资本变 1，才能解锁利润 2 和 3 的项目。

\textbf{暴力基线。} 慢版本每轮扫描全部候选来找最大或最小，再更新候选集合。它的答案选择过程清楚，但动态集合每变化一次就重新找极值，代价太高。放到本题就是：先按“当前资本决定可解锁项目集合，目标是在最多 k 次选择中最大化最终资本”完整试慢版本；样例里已经能看到“规律是按资本排序解锁候选，用利润大顶堆每轮选择当前可做项目里的最大利润”，所以优化前必须先能说清慢版本枚举了哪些候选。

\textbf{暴力过程。} 拿 k=2、w=0、profits=[1, 2, 3]、capital=[0, 1, 1] 手算，初始只能做资本 0 的项目，做完利润 1 后资本变 1，才能解锁利润 2 和 3 的项目。

\textbf{重复工作。} 题面模型：当前资本决定可解锁项目集合，目标是在最多 k 次选择中最大化最终资本。暴力版的慢点要落到本题：慢版本会围绕“当前资本决定可解锁项目集合，目标是在最多 k 次选择中最大化最终资本”反复扫描动态候选集找极值。样例规律是“规律是按资本排序解锁候选，用利润大顶堆每轮选择当前可做项目里的最大利润”，说明每轮真正需要的是堆顶边界，而不是把候选全集重新排序。后面的优化只消掉这类重复工作，不能改变答案集合。

\textbf{优化条件。} 本题的关键观察是：项目按 capital 排序，用指针把可负担项目加入利润大顶堆，每轮选利润最高。这句话必须能翻译成可维护状态；如果题目条件改变后观察不再成立，就不能继续套原来的写法。

\textbf{相邻状态。} 规律是按资本排序解锁候选，用利润大顶堆每轮选择当前可做项目里的最大利润。把这个特例继续拆开看：堆顶必须正好代表下一步最需要处理的候选；若堆里可能有过期对象，读取堆顶前要先清理。堆只解决动态极值，不自动证明被弹出或被丢弃的元素无用。

\textbf{状态复用。} 把每轮全量扫描改成堆顶查询；插入、弹出和过期清理共同维护动态候选集。本题落点是：规律是按资本排序解锁候选，用利润大顶堆每轮选择当前可做项目里的最大利润。手算时只改被新输入影响的那一小部分，而不是回头重算完整候选。

\textbf{指针和字段。} 状态字段要能说清：堆元素字段、堆顶比较规则、候选是否过期、答案读取时机；如果需要懒删除，还要保存删除计数。手算时按这个协议跟踪：拿 k=2、w=0、profits=[1, 2, 3]、capital=[0, 1, 1] 手算，初始只能做资本 0 的项目，做完利润 1 后资本变 1，才能解锁利润 2 和 3 的项目。然后按协议复查：手算时记录堆顶、候选集合和是否有过期元素。每次使用堆顶前都要确认它仍然属于当前问题。

\textbf{代码落点。} 比较器方向先用小样例验证；每次使用堆顶前清理过期项；堆中保存下标时要能回查原数据。关键代码行必须能逐句翻译成“旧状态怎样变成新状态”，否则就只是模板。

\textbf{正确性。} 证明从候选覆盖开始：所有可能成为下一步答案的对象都在堆中，堆顶过期时不会被使用。

\textbf{边界实验。} 先用初始资本为零、没有可做项目、k 大于项目数、利润为零、资本解锁链式增长做反例检查。实验时覆盖初始资本为零、没有可做项目、k 大于项目数、利润为零、资本解锁链式增长，记录堆顶是否过期、比较器是否让正确候选在堆顶、K 或窗口大小为边界值时是否稳定。

\textbf{复杂度变化。} 暴力每轮扫描或排序动态候选，会把线性扫描或排序反复发生；堆把取极值、插入和删除边界控制在对数级。

\textbf{迁移追问。} 如果题目改成“若项目有依赖关系，需要拓扑或约束优化，不能只按资本解锁”，先判断原状态是否仍然覆盖新答案集合，再决定是微调边界、改 value 语义，还是换算法。

\subsection{LeetCode 767 重构字符串}

\textbf{题目说明。} LeetCode 767 重构字符串 的题面入口要先还原成具体任务：相邻不能相同，剩余次数最多的字符最难安排。

\textbf{样例入口。} 拿 s=aab 手算，a 剩 2 次、b 剩 1 次，每轮取两个最高频字符放入答案，可得到 ab 后再放回剩余 a，最后是 aba。拿 aaab 时 a 出现 3 次超过 (4+1)/2，不可能隔开。

\textbf{暴力基线。} 慢版本每轮扫描全部候选来找最大或最小，再更新候选集合。它的答案选择过程清楚，但动态集合每变化一次就重新找极值，代价太高。放到本题就是：先按“相邻不能相同，剩余次数最多的字符最难安排”完整试慢版本；样例里已经能看到“规律是最大频次先判不可能，构造时堆里每轮取两个不同最高频字符”，所以优化前必须先能说清慢版本枚举了哪些候选。

\textbf{暴力过程。} 拿 s=aab 手算，a 剩 2 次、b 剩 1 次，每轮取两个最高频字符放入答案，可得到 ab 后再放回剩余 a，最后是 aba。拿 aaab 时 a 出现 3 次超过 (4+1)/2，不可能隔开。

\textbf{重复工作。} 题面模型：相邻不能相同，剩余次数最多的字符最难安排。暴力版的慢点要落到本题：慢版本会围绕“相邻不能相同，剩余次数最多的字符最难安排”反复扫描动态候选集找极值。样例规律是“规律是最大频次先判不可能，构造时堆里每轮取两个不同最高频字符”，说明每轮真正需要的是堆顶边界，而不是把候选全集重新排序。后面的优化只消掉这类重复工作，不能改变答案集合。

\textbf{优化条件。} 本题的关键观察是：大顶堆每次取两个最高频不同字符放入答案，剩余频次再放回；先检查最大频次是否超过 (n+1)/2。这句话必须能翻译成可维护状态；如果题目条件改变后观察不再成立，就不能继续套原来的写法。

\textbf{相邻状态。} 规律是最大频次先判不可能，构造时堆里每轮取两个不同最高频字符。把这个特例继续拆开看：堆顶必须正好代表下一步最需要处理的候选；若堆里可能有过期对象，读取堆顶前要先清理。堆只解决动态极值，不自动证明被弹出或被丢弃的元素无用。

\textbf{状态复用。} 把每轮全量扫描改成堆顶查询；插入、弹出和过期清理共同维护动态候选集。本题落点是：规律是最大频次先判不可能，构造时堆里每轮取两个不同最高频字符。手算时只改被新输入影响的那一小部分，而不是回头重算完整候选。

\textbf{指针和字段。} 状态字段要能说清：堆元素字段、堆顶比较规则、候选是否过期、答案读取时机；如果需要懒删除，还要保存删除计数。手算时按这个协议跟踪：拿 s=aab 手算，a 剩 2 次、b 剩 1 次，每轮取两个最高频字符放入答案，可得到 ab 后再放回剩余 a，最后是 aba。拿 aaab 时 a 出现 3 次超过 (4+1)/2，不可能隔开。然后按协议复查：手算时记录堆顶、候选集合和是否有过期元素。每次使用堆顶前都要确认它仍然属于当前问题。

\textbf{代码落点。} 比较器方向先用小样例验证；每次使用堆顶前清理过期项；堆中保存下标时要能回查原数据。关键代码行必须能逐句翻译成“旧状态怎样变成新状态”，否则就只是模板。

\textbf{正确性。} 证明从候选覆盖开始：所有可能成为下一步答案的对象都在堆中，堆顶过期时不会被使用。

\textbf{边界实验。} 先用单字符、最高频过半、多个字符同频、最后只剩一个字符、结果不唯一做反例检查。实验时覆盖单字符、最高频过半、多个字符同频、最后只剩一个字符、结果不唯一，记录堆顶是否过期、比较器是否让正确候选在堆顶、K 或窗口大小为边界值时是否稳定。

\textbf{复杂度变化。} 暴力每轮扫描或排序动态候选，会把线性扫描或排序反复发生；堆把取极值、插入和删除边界控制在对数级。

\textbf{迁移追问。} 如果题目改成“若要求相同字符至少间隔 k，堆还要配合冷却队列”，先判断原状态是否仍然覆盖新答案集合，再决定是微调边界、改 value 语义，还是换算法。

\subsection{LeetCode 630 课程表 III}

\textbf{题目说明。} LeetCode 630 课程表 III 的题面入口要先还原成具体任务：每门课有持续时间和截止日，已选课程总时长不能超过当前截止日。

\textbf{样例入口。} 拿课程 [100, 200]、[200, 1300]、[1000, 1250] 手算，按截止日排序后若总时长超出当前截止日，就丢掉已选课程里耗时最长的。

\textbf{暴力基线。} 暴力枚举课程子集和排列，检查是否能在各自截止日前完成，取最多课程数。

\textbf{暴力过程。} 课程 [100, 200]、[200, 1300]、[1000, 1250] 中，按截止日处理时，若加入一门课后总时长超出当前截止日，就要考虑删掉已选中耗时最长的课。

\textbf{重复工作。} 题面模型：每门课有持续时间和截止日，已选课程总时长不能超过当前截止日。暴力版的慢点要落到本题：浪费在枚举所有选择。按截止日排序后，当前前缀课程只需要让总耗时尽量小；超时删除最长课程给未来留下最多时间。后面的优化只消掉这类重复工作，不能改变答案集合。

\textbf{优化条件。} 本题的关键观察是：按截止日排序，先选当前课；若超时，就用大顶堆丢掉已选中耗时最长的课。这句话必须能翻译成可维护状态；如果题目条件改变后观察不再成立，就不能继续套原来的写法。

\textbf{相邻状态。} 规律是大顶堆保存已选课程时长，替换最长课能给未来留下最多时间。把这个特例继续拆开看：先构造一个非贪心选择，再比较两者留下的资源、右边界或剩余时间。只有能把非贪心第一步交换成贪心第一步且不变差，局部选择才是规律。

\textbf{状态复用。} 按 deadline 升序排序。扫描课程时先加入 duration 到 maxHeap 并累加 total\_time；若 total\_time>deadline，就弹出堆中最长 duration。手算时只改被新输入影响的那一小部分，而不是回头重算完整候选。

\textbf{指针和字段。} 状态字段要能说清：total\_time 是已选课程总耗时，maxHeap 保存已选课程时长，deadline 是当前必须满足的最紧约束。手算时按这个协议跟踪：加入 1000 后若超过截止日，从已选课程中丢掉耗时最长的一门，课程数量尽量保留，且总耗时降得最多。

\textbf{代码落点。} 关键是大顶堆保存 duration，不保存收益。每次超时只弹一次最长课，因为刚加入前状态已经满足之前截止日。关键代码行必须能逐句翻译成“旧状态怎样变成新状态”，否则就只是模板。

\textbf{正确性。} 在相同课程数量下，总耗时越短越不妨碍后续课程；超时后删除最长课程能得到同数量减一情况下最短总耗时。

\textbf{边界实验。} 先用课程无法单独完成、截止日相同、替换已选课程、空列表做反例检查。实验时除了课程无法单独完成、截止日相同、替换已选课程、空列表，还要故意替换成一个看似合理但错误的局部规则，构造反例。能解释错误贪心为何失败，才算理解正确贪心。

\textbf{复杂度变化。} 暴力指数级；排序 O(n log n)，每门课入堆一次、可能出堆一次，O(n log n) 时间，O(n) 空间。

\textbf{迁移追问。} 如果题目改成“这是反悔贪心：接受后在约束被破坏时删最差候选”，先判断原状态是否仍然覆盖新答案集合，再决定是微调边界、改 value 语义，还是换算法。

\subsection{LeetCode 435 无重叠区间}

\textbf{题目说明。} LeetCode 435 无重叠区间 的题面入口要先还原成具体任务：保留最多不重叠区间等价于删除最少区间。

\textbf{样例入口。} 拿 [[1, 2], [2, 3], [3, 4], [1, 3]] 手算，选 [1, 2] 后还能接 [2, 3] 和 [3, 4]；若先选 [1, 3]，会挡住更多区间。

\textbf{暴力基线。} 暴力可以枚举保留哪些区间，再检查它们是否互不重叠，最后用总数减最大保留数量。

\textbf{暴力过程。} 以 [[1, 2], [2, 3], [3, 4], [1, 3]] 为例，选 [1, 2] 后还能接 [2, 3] 和 [3, 4]；若先选 [1, 3] 会挡住更多后续区间。

\textbf{重复工作。} 题面模型：保留最多不重叠区间等价于删除最少区间。暴力版的慢点要落到本题：浪费在枚举保留集合。区间选择里结束越早，给后面留下的空间越大，可以作为安全局部选择。后面的优化只消掉这类重复工作，不能改变答案集合。

\textbf{优化条件。} 本题的关键观察是：按终点升序选择，结束越早给后续留下空间越多。这句话必须能翻译成可维护状态；如果题目条件改变后观察不再成立，就不能继续套原来的写法。

\textbf{相邻状态。} 规律是按右端点升序选择，结束越早给未来留下的空间越多，删除数等于总数减保留数。把这个特例继续拆开看：先构造一个非贪心选择，再比较两者留下的资源、右边界或剩余时间。只有能把非贪心第一步交换成贪心第一步且不变差，局部选择才是规律。

\textbf{状态复用。} 按右端点升序排序。维护 last\_end；若当前区间 start>=last\_end，就保留它并更新 last\_end，否则删除它。手算时只改被新输入影响的那一小部分，而不是回头重算完整候选。

\textbf{指针和字段。} 状态字段要能说清：last\_end 是已保留区间的最后右端，kept 是保留数量，removed 是删除数量。手算时按这个协议跟踪：排序后先保留 [1, 2]；[1, 3] 与它重叠删除；[2, 3] 和 [3, 4] 都能接上。

\textbf{代码落点。} 关键是端点相接是否算重叠。本题中 [1, 2] 和 [2, 3] 不重叠，所以条件是 start>=last\_end。关键代码行必须能逐句翻译成“旧状态怎样变成新状态”，否则就只是模板。

\textbf{正确性。} 任意最优解的第一个区间都可替换为结束最早的区间，后续可用空间不变小；重复此过程得到贪心解。

\textbf{边界实验。} 先用端点相接是否重叠、完全覆盖、相同终点、空列表做反例检查。实验时除了端点相接是否重叠、完全覆盖、相同终点、空列表，还要故意替换成一个看似合理但错误的局部规则，构造反例。能解释错误贪心为何失败，才算理解正确贪心。

\textbf{复杂度变化。} 暴力指数级；排序 O(n log n)，扫描 O(n)，空间取决于排序。

\textbf{迁移追问。} 如果题目改成“用交换论证说明最早结束的选择不劣于其他第一选择”，先判断原状态是否仍然覆盖新答案集合，再决定是微调边界、改 value 语义，还是换算法。

\subsection{LeetCode 406 根据身高重建队列}

\textbf{题目说明。} LeetCode 406 根据身高重建队列 的题面入口要先还原成具体任务：高个子的位置先确定后，不会被矮个子影响其前方高个数量。

\textbf{样例入口。} 拿 [7, 0]、[7, 1]、[5, 0] 手算，先放高个子，矮个子插入不会改变高个子前面高个数量；[5, 0] 插到最前也不影响 [7, 0] 的统计。

\textbf{暴力基线。} 慢版本枚举选择顺序或用 DP 保留多个可能方案，先确认最优值存在。随后才检查局部选择是否能通过交换或领先性固定下来。放到本题就是：先按“高个子的位置先确定后，不会被矮个子影响其前方高个数量”完整试慢版本；样例里已经能看到“规律是按身高降序、k 升序排序，再按 k 插入”，所以优化前必须先能说清慢版本枚举了哪些候选。

\textbf{暴力过程。} 拿 [7, 0]、[7, 1]、[5, 0] 手算，先放高个子，矮个子插入不会改变高个子前面高个数量；[5, 0] 插到最前也不影响 [7, 0] 的统计。

\textbf{重复工作。} 题面模型：高个子的位置先确定后，不会被矮个子影响其前方高个数量。暴力版的慢点要落到本题：慢版本会围绕“高个子的位置先确定后，不会被矮个子影响其前方高个数量”枚举选择顺序。样例规律是“规律是按身高降序、k 升序排序，再按 k 插入”，说明存在一个局部选择或边界状态可以安全保留；不能证明这个选择不变差时，就不能把它叫贪心。后面的优化只消掉这类重复工作，不能改变答案集合。

\textbf{优化条件。} 本题的关键观察是：按身高降序、k 升序排序，再按 k 插入当前结果位置。这句话必须能翻译成可维护状态；如果题目条件改变后观察不再成立，就不能继续套原来的写法。

\textbf{相邻状态。} 规律是按身高降序、k 升序排序，再按 k 插入。把这个特例继续拆开看：先构造一个非贪心选择，再比较两者留下的资源、右边界或剩余时间。只有能把非贪心第一步交换成贪心第一步且不变差，局部选择才是规律。

\textbf{状态复用。} 把指数级选择压成一个可证明安全的局部选择；状态通常是当前最远边界、最早结束时间、最小代价或剩余资源。本题落点是：规律是按身高降序、k 升序排序，再按 k 插入。手算时只改被新输入影响的那一小部分，而不是回头重算完整候选。

\textbf{指针和字段。} 状态字段要能说清：处理顺序、当前已选方案、剩余资源或最远边界、答案计数；若使用排序，要说明排序键；每次局部选择后要能说明当前状态不落后。手算时按这个协议跟踪：拿 [7, 0]、[7, 1]、[5, 0] 手算，先放高个子，矮个子插入不会改变高个子前面高个数量；[5, 0] 插到最前也不影响 [7, 0] 的统计。然后按协议复查：手算时同时写一个非贪心选择，与贪心选择比较剩余资源。若无法说明贪心后剩余资源不差，就不能直接下结论。

\textbf{代码落点。} 把局部规则写成业务含义；若需要排序，就处理相同关键字；循环中只维护证明需要的状态，不把局部选择和答案更新混在一起。关键代码行必须能逐句翻译成“旧状态怎样变成新状态”，否则就只是模板。

\textbf{正确性。} 证明从交换或领先性开始：任意最优解都能替换成当前贪心选择而不变差，或贪心状态始终不落后。

\textbf{边界实验。} 先用相同身高、k 为零、插入到末尾、结果稳定性做反例检查。实验时除了相同身高、k 为零、插入到末尾、结果稳定性，还要故意替换成一个看似合理但错误的局部规则，构造反例。能解释错误贪心为何失败，才算理解正确贪心。

\textbf{复杂度变化。} 暴力枚举选择顺序可能指数级；贪心通常先排序，再线性扫描或用堆维护少量候选。

\textbf{迁移追问。} 如果题目改成“这是排序后插入的构造型贪心，证明依赖高个子先固定”，先判断原状态是否仍然覆盖新答案集合，再决定是微调边界、改 value 语义，还是换算法。

\subsection{LeetCode 860 柠檬水找零}

\textbf{题目说明。} LeetCode 860 柠檬水找零 的题面入口要先还原成具体任务：找零资源只有 5 元和 10 元，5 元比 10 元更灵活。

\textbf{样例入口。} 拿 bills=[5, 5, 5, 10, 20] 手算，前三个人给 5 后手里有三个 5；收到 10 找一个 5；收到 20 时优先找 10+5，保留更多 5 给后续 10 找零。

\textbf{暴力基线。} 慢版本枚举选择顺序或用 DP 保留多个可能方案，先确认最优值存在。随后才检查局部选择是否能通过交换或领先性固定下来。放到本题就是：先按“找零资源只有 5 元和 10 元，5 元比 10 元更灵活”完整试慢版本；样例里已经能看到“规律是维护 5 元和 10 元数量，20 找零时优先使用较不灵活的 10 元”，所以优化前必须先能说清慢版本枚举了哪些候选。

\textbf{暴力过程。} 拿 bills=[5, 5, 5, 10, 20] 手算，前三个人给 5 后手里有三个 5；收到 10 找一个 5；收到 20 时优先找 10+5，保留更多 5 给后续 10 找零。

\textbf{重复工作。} 题面模型：找零资源只有 5 元和 10 元，5 元比 10 元更灵活。暴力版的慢点要落到本题：慢版本会围绕“找零资源只有 5 元和 10 元，5 元比 10 元更灵活”枚举选择顺序。样例规律是“规律是维护 5 元和 10 元数量，20 找零时优先使用较不灵活的 10 元”，说明存在一个局部选择或边界状态可以安全保留；不能证明这个选择不变差时，就不能把它叫贪心。后面的优化只消掉这类重复工作，不能改变答案集合。

\textbf{优化条件。} 本题的关键观察是：收到 10 用一个 5 找零；收到 20 优先用 10+5，不能时再用三个 5。这句话必须能翻译成可维护状态；如果题目条件改变后观察不再成立，就不能继续套原来的写法。

\textbf{相邻状态。} 规律是维护 5 元和 10 元数量，20 找零时优先使用较不灵活的 10 元。把这个特例继续拆开看：先构造一个非贪心选择，再比较两者留下的资源、右边界或剩余时间。只有能把非贪心第一步交换成贪心第一步且不变差，局部选择才是规律。

\textbf{状态复用。} 把指数级选择压成一个可证明安全的局部选择；状态通常是当前最远边界、最早结束时间、最小代价或剩余资源。本题落点是：规律是维护 5 元和 10 元数量，20 找零时优先使用较不灵活的 10 元。手算时只改被新输入影响的那一小部分，而不是回头重算完整候选。

\textbf{指针和字段。} 状态字段要能说清：处理顺序、当前已选方案、剩余资源或最远边界、答案计数；若使用排序，要说明排序键；每次局部选择后要能说明当前状态不落后。手算时按这个协议跟踪：拿 bills=[5, 5, 5, 10, 20] 手算，前三个人给 5 后手里有三个 5；收到 10 找一个 5；收到 20 时优先找 10+5，保留更多 5 给后续 10 找零。然后按协议复查：手算时同时写一个非贪心选择，与贪心选择比较剩余资源。若无法说明贪心后剩余资源不差，就不能直接下结论。

\textbf{代码落点。} 把局部规则写成业务含义；若需要排序，就处理相同关键字；循环中只维护证明需要的状态，不把局部选择和答案更新混在一起。关键代码行必须能逐句翻译成“旧状态怎样变成新状态”，否则就只是模板。

\textbf{正确性。} 证明从交换或领先性开始：任意最优解都能替换成当前贪心选择而不变差，或贪心状态始终不落后。

\textbf{边界实验。} 先用第一张不是 5、连续 10、连续 20、只有 5、5 元数量不足做反例检查。实验时除了第一张不是 5、连续 10、连续 20、只有 5、5 元数量不足，还要故意替换成一个看似合理但错误的局部规则，构造反例。能解释错误贪心为何失败，才算理解正确贪心。

\textbf{复杂度变化。} 暴力枚举选择顺序可能指数级；贪心通常先排序，再线性扫描或用堆维护少量候选。

\textbf{迁移追问。} 如果题目改成“若面额集合变化，需要重新证明哪种资源更灵活”，先判断原状态是否仍然覆盖新答案集合，再决定是微调边界、改 value 语义，还是换算法。

\subsection{LeetCode 871 最低加油次数}

\textbf{题目说明。} LeetCode 871 最低加油次数 的题面入口要先还原成具体任务：行驶过程中，经过的加油站形成可反悔的燃料候选。

\textbf{样例入口。} 拿 target=100、startFuel=10 手算，车先到 10 号站，把 60 升油加入最大堆；之后油量不足时，从已经路过的站里取当前最大油量。

\textbf{暴力基线。} 慢版本每轮扫描全部候选来找最大或最小，再更新候选集合。它的答案选择过程清楚，但动态集合每变化一次就重新找极值，代价太高。放到本题就是：先按“行驶过程中，经过的加油站形成可反悔的燃料候选”完整试慢版本；样例里已经能看到“规律是能到的站先入堆，需要加油时再反悔选择最大燃料，保证次数最少”，所以优化前必须先能说清慢版本枚举了哪些候选。

\textbf{暴力过程。} 拿 target=100、startFuel=10 手算，车先到 10 号站，把 60 升油加入最大堆；之后油量不足时，从已经路过的站里取当前最大油量。

\textbf{重复工作。} 题面模型：行驶过程中，经过的加油站形成可反悔的燃料候选。暴力版的慢点要落到本题：慢版本会围绕“行驶过程中，经过的加油站形成可反悔的燃料候选”反复扫描动态候选集找极值。样例规律是“规律是能到的站先入堆，需要加油时再反悔选择最大燃料，保证次数最少”，说明每轮真正需要的是堆顶边界，而不是把候选全集重新排序。后面的优化只消掉这类重复工作，不能改变答案集合。

\textbf{优化条件。} 本题的关键观察是：按位置扫描，油量不足到达下一点时，从已路过站点中取最大燃料补充。这句话必须能翻译成可维护状态；如果题目条件改变后观察不再成立，就不能继续套原来的写法。

\textbf{相邻状态。} 规律是能到的站先入堆，需要加油时再反悔选择最大燃料，保证次数最少。把这个特例继续拆开看：堆顶必须正好代表下一步最需要处理的候选；若堆里可能有过期对象，读取堆顶前要先清理。堆只解决动态极值，不自动证明被弹出或被丢弃的元素无用。

\textbf{状态复用。} 把每轮全量扫描改成堆顶查询；插入、弹出和过期清理共同维护动态候选集。本题落点是：规律是能到的站先入堆，需要加油时再反悔选择最大燃料，保证次数最少。手算时只改被新输入影响的那一小部分，而不是回头重算完整候选。

\textbf{指针和字段。} 状态字段要能说清：堆元素字段、堆顶比较规则、候选是否过期、答案读取时机；如果需要懒删除，还要保存删除计数。手算时按这个协议跟踪：拿 target=100、startFuel=10 手算，车先到 10 号站，把 60 升油加入最大堆；之后油量不足时，从已经路过的站里取当前最大油量。然后按协议复查：手算时记录堆顶、候选集合和是否有过期元素。每次使用堆顶前都要确认它仍然属于当前问题。

\textbf{代码落点。} 比较器方向先用小样例验证；每次使用堆顶前清理过期项；堆中保存下标时要能回查原数据。关键代码行必须能逐句翻译成“旧状态怎样变成新状态”，否则就只是模板。

\textbf{正确性。} 证明从候选覆盖开始：所有可能成为下一步答案的对象都在堆中，堆顶过期时不会被使用。

\textbf{边界实验。} 先用起点油足够、无法到达第一个站、多个站同位置、目标作为终点事件做反例检查。实验时覆盖起点油足够、无法到达第一个站、多个站同位置、目标作为终点事件，记录堆顶是否过期、比较器是否让正确候选在堆顶、K 或窗口大小为边界值时是否稳定。

\textbf{复杂度变化。} 暴力每轮扫描或排序动态候选，会把线性扫描或排序反复发生；堆把取极值、插入和删除边界控制在对数级。

\textbf{迁移追问。} 如果题目改成“这是反悔贪心，堆里保存的是已经经过但尚未使用的资源”，先判断原状态是否仍然覆盖新答案集合，再决定是微调边界、改 value 语义，还是换算法。

\subsection{LeetCode 215 数组中的第 K 个最大元素}

\textbf{题目说明。} LeetCode 215 数组中的第 K 个最大元素 的题面入口要先还原成具体任务：只关心第 K 大，不需要完整排序。

\textbf{样例入口。} 拿 [3, 2, 1, 5, 6, 4]、k=2 手算，固定大小为 2 的小顶堆最终保存最大的两个数，堆顶就是第 2 大。

\textbf{暴力基线。} 暴力排序整个数组，然后返回降序第 k 个或升序第 n-k 个。

\textbf{暴力过程。} 以 [3, 2, 1, 5, 6, 4]、k=2 为例，完整排序后第 2 大是 5；但为了一个第 K 大，不需要知道所有元素的完整顺序。

\textbf{重复工作。} 题面模型：只关心第 K 大，不需要完整排序。暴力版的慢点要落到本题：浪费在完整排序。我们只需要维护当前最大的 K 个元素，其中最小的那个就是第 K 大边界。后面的优化只消掉这类重复工作，不能改变答案集合。

\textbf{优化条件。} 本题的关键观察是：大小为 K 的小顶堆保存当前前 K 大边界；快速选择可做到平均线性。这句话必须能翻译成可维护状态；如果题目条件改变后观察不再成立，就不能继续套原来的写法。

\textbf{相邻状态。} 规律是堆里只保留当前前 k 大候选，比堆顶小的数不可能进入答案。把这个特例继续拆开看：堆顶必须正好代表下一步最需要处理的候选；若堆里可能有过期对象，读取堆顶前要先清理。堆只解决动态极值，不自动证明被弹出或被丢弃的元素无用。

\textbf{状态复用。} 大小为 K 的小顶堆保存当前前 K 大。新数大于堆顶时，弹出堆顶并加入新数；小于等于堆顶则不可能进入前 K。手算时只改被新输入影响的那一小部分，而不是回头重算完整候选。

\textbf{指针和字段。} 状态字段要能说清：heap.top 是当前前 K 大中的最小值，size 保持不超过 K，num 是当前扫描元素。手算时按这个协议跟踪：扫描样例时，堆最终保存 5 和 6，堆顶 5 就是第 2 大。

\textbf{代码落点。} 关键是使用小顶堆，C++ priority\_queue 默认大顶堆，需要 greater<int>。若选择快速选择，要处理随机 pivot 和重复值。关键代码行必须能逐句翻译成“旧状态怎样变成新状态”，否则就只是模板。

\textbf{正确性。} 堆中始终保存已扫描元素的前 K 大；若新元素不超过堆顶，它最多排在第 K 大之后，可以安全丢弃。

\textbf{边界实验。} 先用K 为一、K 为 n、重复值、随机 pivot做反例检查。实验时覆盖K 为一、K 为 n、重复值、随机 pivot，记录堆顶是否过期、比较器是否让正确候选在堆顶、K 或窗口大小为边界值时是否稳定。

\textbf{复杂度变化。} 排序 O(n log n)，小顶堆 O(n log k) 时间、O(k) 空间；快速选择平均 O(n)。

\textbf{迁移追问。} 如果题目改成“若数据流持续到来，堆方案比一次性快速选择更自然”，先判断原状态是否仍然覆盖新答案集合，再决定是微调边界、改 value 语义，还是换算法。

\begin{principlebox}{本节复盘方式}
不要只看最终算法名。每道题复盘时先遮住优化版，口述暴力枚举对象；再指出相邻窗口、历史前缀、候选区间、搜索状态或子问题里哪一部分被重复处理；最后用一个边界样例验证状态更新顺序。能讲完这三步，才算真正掌握本章案例。
\end{principlebox}
